% =====================================================================
% Sekcja 5 - Kalibracja
% Wartości z overnight_results/RESULTS_CALIBRATION_15B.md (calibration_15b.png).
% =====================================================================

\section{Kalibracja}
\label{sec:kalibracja}

Wyniki z sekcji~\ref{sec:wyniki} pokazują, że pewność modelu
koreluje z poprawnością. Korelacja jest jednak własnością czysto
porządkową - mówi tylko tyle, że pytania da się ustawić od
najbardziej do najmniej ryzykownego. Nie odpowiada natomiast na
pytanie, czy pewność rzędu 0,85 faktycznie oznacza 85\% szans na
poprawną odpowiedź. To już kwestia kalibracji i właśnie jej
dotyczy ta sekcja.

Jako pewność modelu przyjęto $1/\mathrm{PPL}$ - liczbę z przedziału
$(0, 1]$. Jej zgodność z poprawnością mierzy Expected Calibration
Error (ECE) na binach o równej liczbie obserwacji
($n_{\text{bins}}=10$). Biny o równej liczności są tu potrzebne,
bo przy dekodowaniu zachłannym wartości $1/\mathrm{PPL}$ skupiają
się blisko 1,0 (sekcja~\ref{sec:dekodowanie}) i biny o stałej
szerokości byłyby w większości puste.

\begin{table}[h]
\centering
\small
\begin{tabular}{lrrr}
\toprule
\textbf{Ustawienie} & $n$ & \textbf{dokładność} & \textbf{ECE} \\
\midrule
1,5B CoT                       & 60 & 60,0\% & $0{,}265$ \\
1,5B Agent (cała traj.)        & 60 & 28,3\% & $0{,}425$ \\
1,5B Agent (segmenty action)   & 48 & 25,0\% & $0{,}445$ \\
\bottomrule
\end{tabular}
\caption{Expected Calibration Error dla pewności zdefiniowanej
jako $1/\mathrm{PPL}$ (binowanie o równej liczności,
$n_{\text{bins}}=10$). Niższe ECE oznacza lepszą zgodność
deklarowanej pewności z faktyczną poprawnością.}
\label{tab:ece}
\end{table}

ECE jest wysokie w każdym ustawieniu (tabela~\ref{tab:ece}).
Diagramy niezawodności (plik \texttt{calibration\_15b.png})
tłumaczą skąd to się bierze: w prawie każdym binie średnia pewność
leży powyżej średniej poprawności, czyli model jest przesadnie
pewny. W wariancie agentowym widać to najmocniej - w najniższych
binach pewności model deklaruje 0,49-0,66, podczas gdy poprawność
w tych binach wynosi 0\%. Nawet tam, gdzie agent jest najmniej
pewny, i tak jest pewniejszy, niż uzasadniałyby to jego wyniki.

Zestawienie tej obserwacji z sekcją~\ref{sec:wyniki} prowadzi do
głównego wniosku. Perplexity liczone na rozumowaniu modelu istotnie
koreluje z poprawnością, a mimo to daje wysokie ECE. Dobry porządek
nie pociąga więc za sobą dobrej kalibracji - jedno z drugiego nie
wynika.

Ma to konkretne znaczenie dla zastosowań. Tam, gdzie liczy się
sama kolejność, perplexity się sprawdza: można nim wybrać
najpewniejszą z kilku odpowiedzi albo wskazać najbardziej ryzykowne
pytania do ręcznego sprawdzenia. Gorzej z regułami opartymi na
stałym progu, w stylu ,,odrocz odpowiedź, gdy pewność spadnie
poniżej 0,7'' - wartość 0,7 nie oznacza tu 70\% szans na poprawność,
lecz znacznie mniej, więc taki próg trzeba najpierw wykalibrować
na danych z konkretnego zadania. Jest to zastrzeżenie wobec prac
omawianych w sekcji~\ref{sec:agent-uncertainty}: kierunek sygnału
jest dobry, ale jego skala nie nadaje się do bezpośredniego progowania.
